\SourcePage{174}{177}
\section{The secular equation}

We now turn to the case in which the unperturbed operator $\hat H_0$ has
degenerate eigenvalues. Let $\psi_n^{(0)}$, $\psi_{n'}^{(0)}$,~\ldots denote
eigenfunctions belonging to the same energy eigenvalue $E_n^{(0)}$. As we
know, the choice of these functions is not unique: they may be replaced by
any $s$ independent linear combinations of the same functions, where $s$ is
the degeneracy of the level $E_n^{(0)}$. This freedom disappears, however,
when the wave functions are required to change only slightly under the
applied small perturbation.
\SourcePage{175}{178}
For the moment, let $\psi_n^{(0)}$, $\psi_{n'}^{(0)}$,~\ldots denote
arbitrarily chosen unperturbed eigenfunctions. The correct zeroth-order
functions are linear combinations of the form
\[
  c_n^{(0)}\psi_n^{(0)}+c_{n'}^{(0)}\psi_{n'}^{(0)}+\ldots
\]
The coefficients in these combinations and the first-order corrections to
the eigenvalues are determined as follows.

Write equations (38.4) for $k=n,n',\ldots$, substituting in first order
$E=E_n^{(0)}+E^{(1)}$. For the $c_k$ it is sufficient to retain their
zeroth-order values $c_n=c_n^{(0)}$, $c_{n'}=c_{n'}^{(0)},\ldots$, while
$c_m=0$ for $m\ne n,n',\ldots$. We then obtain
\[
  E^{(1)}c_n^{(0)}=\sum_{n'}V_{nn'}c_{n'}^{(0)},
\]
or
\begin{equation}
  \sum_{n'}(V_{nn'}-E^{(1)}\delta_{nn'})c_{n'}^{(0)}=0,
  \tag{39.1}
\end{equation}
where $n$ and $n'$ run over all values labeling states that belong to the
given unperturbed eigenvalue $E_n^{(0)}$. This homogeneous system of linear
equations for the $c_n^{(0)}$ has nonzero solutions if the determinant of its
coefficients vanishes. Thus
\begin{equation}
  |V_{nn'}-E^{(1)}\delta_{nn'}|=0.
  \tag{39.2}
\end{equation}
This equation has degree $s$ in $E^{(1)}$ and, in general, has $s$ distinct
real roots. These roots are the required first-order corrections to the
eigenvalues. Equation (39.2) is called the \emph{secular
equation}\footnote{Also called the age equation (the French word
\emph{siècle} means an age); the name was borrowed from celestial
mechanics.}. The sum of its roots is the sum of the diagonal matrix elements
$V_{nn}$, $V_{n'n'},\ldots$; it is the coefficient of $E^{(1)s-1}$ in the
equation.

Substituting the roots of (39.2) one by one into (39.1) and solving the
resulting systems gives the coefficients $c_n^{(0)}$ and hence the correct
zeroth-order eigenfunctions.
\SourcePage{176}{179}
As a result of the perturbation, the initially degenerate energy level is in
general no longer degenerate, since the roots of (39.2) are in general
different. The perturbation is said to lift the degeneracy. The lifting may
be complete or partial; in the latter case the perturbed level remains
degenerate, but with a lower multiplicity than before.

It may happen that, for some reason, all matrix elements for transitions
within one group of mutually degenerate states $n,n',\ldots$ are especially
small, or even vanish. It may then be useful, while retaining $V_{nn'}$ in
first order, also to include in higher orders the matrix elements $V_{nm}$,
with $m\ne n,n',\ldots$, for transitions to states of other energies. We
shall do this by including $V_{mn}$ through second order.

In (38.4) with $k=n$, set $E=E_n^{(0)}+E^{(1)}$ on the left-hand side; we
retain the notation $E^{(1)}$ for the energy correction in the approximation
now considered. Replace $c_n$ by $c_n^{(0)}$. Since $c_m^{(0)}=0$ for every
$m\ne n,n',\ldots$, we have
\begin{equation}
  E^{(1)}c_n^{(0)}=\sum_mV_{nm}c_m^{(1)}
  +\sum_{n'}V_{nn'}c_{n'}^{(0)}.
  \tag{39.3}
\end{equation}
Equations (38.4) with $k=m\ne n,n',\ldots$, through first order, give
\[
  (E_n^{(0)}-E_m^{(0)})c_m^{(1)}
  =\sum_{n'}V_{mn'}c_{n'}^{(0)},
\]
and hence
\[
  c_m^{(1)}=\sum_{n'}\frac{V_{mn'}}{E_n^{(0)}-E_m^{(0)}}c_{n'}^{(0)}.
\]
Substitution in (39.3) gives
\[
  E^{(1)}c_n^{(0)}=\sum_{n'}c_{n'}^{(0)}
  \left(V_{nn'}+\sum_m\frac{V_{nm}V_{mn'}}{E_n^{(0)}-E_m^{(0)}}\right).
\]
This system now replaces (39.1). Its compatibility condition again gives a
secular equation, differing from (39.2) by the replacement
\begin{equation}
  V_{nn'}\longrightarrow V_{nn'}
  +\sum_m\frac{V_{nm}V_{mn'}}{E_n^{(0)}-E_m^{(0)}}.
  \tag{39.4}
\end{equation}
\SourcePage{177}{180}
\ProblemHeading
\textbf{1.} Find the first-order corrections to the eigenvalue and the
correct zeroth-order functions for a doubly degenerate level.

\SolutionHeading Equation (39.2) now has the form
\[
  \begin{vmatrix}
   V_{11}-E^{(1)}&V_{21}\\
   V_{12}&V_{22}-E^{(1)}
  \end{vmatrix}=0
\]
(the indices 1 and 2 refer to two arbitrarily chosen unperturbed
eigenfunctions $\psi_1^{(0)}$ and $\psi_2^{(0)}$ of the degenerate level).
Solving it, we find
\begin{equation}
  E^{(1)}=\frac12[V_{11}+V_{22}\pm\hbar\omega^{(1)}],
  \tag{1}
\end{equation}
where
\[
  \hbar\omega^{(1)}=\sqrt{(V_{11}-V_{22})^2+4|V_{12}|^2}
\]
denotes the difference between the two values of the correction $E^{(1)}$.
Solving (39.1) with these values of $E^{(1)}$, we obtain the coefficients in
the normalized correct zeroth-order functions
$\psi^{(0)}=c_1^{(0)}\psi_1^{(0)}+c_2^{(0)}\psi_2^{(0)}$:
\begin{equation}
 \begin{aligned}
  c_1^{(0)}&=\left\{\frac{V_{12}}{2|V_{12}|}
    \left[1\pm\frac{V_{11}-V_{22}}{\hbar\omega^{(1)}}\right]\right\}^{1/2},\\
  c_2^{(0)}&=\pm\left\{\frac{V_{21}}{2|V_{12}|}
    \left[1\mp\frac{V_{11}-V_{22}}{\hbar\omega^{(1)}}\right]\right\}^{1/2}.
 \end{aligned}
 \tag{2}
\end{equation}

\ProblemHeading
\textbf{2.} Derive formulas for the first-order corrections to the
eigenfunctions and the second-order corrections to the eigenvalues.

\SolutionHeading Let the functions $\psi_n^{(0)}$ be chosen as the correct
zeroth-order functions. The matrix $V_{nn'}$ defined with these functions is
evidently diagonal in the indices $n,n'$ belonging to the same group of
functions of a degenerate level. Its diagonal elements $V_{nn}$,
$V_{n'n'},\ldots$ equal the corresponding first-order corrections
$E_n^{(1)}$, $E_{n'}^{(1)},\ldots$.

Consider the perturbation of the eigenfunction $\psi_n^{(0)}$. In zeroth
order, $E=E_n^{(0)}$, $c_n^{(0)}=1$, and $c_m^{(0)}=0$ for $m\ne n$. In
first order, $E=E_n^{(0)}+V_{nn}$, $c_n=1+c_n^{(1)}$, and
$c_m=c_m^{(1)}$. From the general system (38.4), take the equation with
$k\ne n,n',\ldots$ and retain its first-order terms:
\[
  (E_n^{(0)}-E_k^{(0)})c_k^{(1)}=V_{kn}c_n^{(0)}=V_{kn},
\]
whence
\begin{equation}
  c_k^{(1)}=\frac{V_{kn}}{E_n^{(0)}-E_k^{(0)}}
  \quad\text{for}\quad k\ne n,n',\ldots
  \tag{1}
\end{equation}
Next write the equation with $k=n'$, retaining its second-order terms:
\[
  E_n^{(1)}c_{n'}^{(1)}=V_{n'n'}c_{n'}^{(1)}
  +\sum_m'V_{n'm}c_m^{(1)}.
\]
\SourcePage{178}{181}
In the sum over $m$, the terms with $m=n,n',\ldots$ are omitted. Substituting
$E_n^{(1)}=V_{nn}$ and (1) for $c_m^{(1)}$, we obtain, for $n'\ne n$,
\begin{equation}
  c_{n'}^{(1)}=\frac1{V_{nn}-V_{n'n'}}
  \sum_m'\frac{V_{n'm}V_{mn}}{E_n^{(0)}-E_m^{(0)}}.
  \tag{2}
\end{equation}
The coefficient $c_n^{(1)}$ is zero in this approximation. Formulas (1) and
(2) determine the first-order correction
$\psi_n^{(1)}=\sum_m c_m^{(1)}\psi_m^{(0)}$ to the
eigenfunctions\footnote{Notice that the requirement that the quantities (1)
and (2) be small, and hence the condition for applicability of this
perturbation method, still requires (38.9) only for transitions between
states belonging to different energy levels. Transitions between states of
the same degenerate level are accounted for by the secular equation in a
certain sense exactly.}.

Finally, retaining the second-order terms in (38.4) with $k=n$, we obtain
the second-order energy correction
\begin{equation}
  E_n^{(2)}=\sum_m'\frac{V_{nm}V_{mn}}{E_n^{(0)}-E_m^{(0)}},
  \tag{3}
\end{equation}
which has the same formal appearance as (38.10).

\ProblemHeading
\textbf{3.} At the initial time $t=0$, a system is in a state
$\psi_1^{(0)}$ belonging to a doubly degenerate level. Find the probability
that at a later time $t$ the system is in the other state $\psi_2^{(0)}$ of
the same energy, when the transition is caused by a constant perturbation.

\SolutionHeading Construct the correct zeroth-order functions:
\[
  \psi=c_1\psi_1+c_2\psi_2,\qquad
  \psi'=c_1'\psi_1+c_2'\psi_2,
\]
where $c_1,c_2$ and $c_1',c_2'$ are the two pairs of coefficients given by
(2) of Problem 1. Superscripts $(0)$ have been omitted from all quantities
for brevity. Conversely,
\[
  \psi_1=\frac{c_2'\psi-c_2\psi'}{c_1c_2'-c_1'c_2}.
\]
The functions $\psi$ and $\psi'$ belong to states with perturbed energies
$E+E^{(1)}$ and $E+E^{(1)'}$, where $E^{(1)}$ and $E^{(1)'}$ are the two
values of correction (1) in Problem 1. Introducing the time factors gives
the time-dependent wave function
\[
  \Psi_1=\frac{\exp[-(i/\hbar)Et]}{c_1c_2'-c_1'c_2}
  \left[c_2'\psi\exp\left(-\frac i\hbar E^{(1)}t\right)
       -c_2\psi'\exp\left(-\frac i\hbar E^{(1)'}t\right)\right].
\]
At $t=0$, $\Psi_1=\psi_1$. Expressing $\psi$ and $\psi'$ once more in terms
of $\psi_1$ and $\psi_2$, we write $\Psi_1$ as a linear combination of
$\psi_1$ and $\psi_2$ with time-dependent coefficients. The squared modulus
of the coefficient of $\psi_2$ is the required transition probability
$w_{21}$. Using (1) and (2) of Problem 1 gives
\[
  w_{21}=2\frac{|V_{12}|^2}{(\hbar\omega^{(1)})^2}
  [1-\cos(\omega^{(1)}t)].
\]
Thus the probability oscillates periodically with frequency
$\omega^{(1)}$. For times $t$ small compared with the corresponding period,
the expression in
\SourcePage{179}{182}
square brackets, and hence $w_{21}$ itself, is proportional to $t^2$:
\[
  w_{21}=\frac1{\hbar^2}|V_{12}|^2t^2.
\]
This formula follows very simply by the method developed in the next section,
using (40.4).
